\subsection{Estrutura do Trabalho}

\par Além do capítulo 1 de introdução, o trabalho está organizado em outros 5 capítulos. No capitulo 2 são abordados os conceitos e tecnologias utilizadas dentro do trabalho, conceitos de automação industrial, tecnologia da informação industrial e desenvolvimento de software e como todos esses conceitos se relacionam para a solução proposta.

\par No Capítulo 3, são especificados os requisitos funcionais e não funcionais e a arquitetura do sistema, com destaque para a separação entre camada de apresentação (web), camada de serviços e camada de obtenção ou persistência dos dados de monitoramento.

\par No Capítulo 4, são implementados os serviços de backend responsáveis pela leitura periódica ou sob demanda, pelo registro e pela disponibilização dos dados de processo de forma modular e interoperável.


\par No Capítulo 5, são implementadas a interface web de monitoramento, permitindo a visualização dos principais indicadores, tendências ou telas resumo associadas ao processo acompanhado.

\par No Capítulo 6, são validados o módulo de monitoramento por meio de testes e de um cenário demonstrativo, evidenciando o correto acesso às variáveis e a estabilidade da solução frente ao volume de dados esperado.